% ============================================================================
%  presentation/scriptstyle.tex --- shared preamble for the speaker scripts
%
%  script.tex (the full one) and script-short.tex (30 seconds a slide) are the
%  same document at two lengths. They looked the same by coincidence until this
%  file existed; now they look the same by construction, and a change to the
%  two-column layout cannot land in one and not the other.
%
%  \scriptname{...} sets the right-hand running head, which is the only thing
%  the two documents disagree about.
% ============================================================================
% Both \scriptname and its default live inside \makeatletter: defining the
% setter outside it made LaTeX read \@scriptname as \@ followed by the letters
% "scriptname", so the setter silently did nothing and both documents kept the
% default running head.
\makeatletter
\newcommand{\scriptname}[1]{\gdef\@scriptname{#1}}
\gdef\@scriptname{EN spoken \; $\cdot$ \; FR pour comprendre}
\makeatother

\usepackage{iftex}
\ifLuaTeX
  \usepackage{fontspec}
  \IfFontExistsTF{Lato}{\setsansfont{Lato}}{}
\else
  \usepackage[utf8]{inputenc}
  \usepackage[T1]{fontenc}
  \usepackage{lmodern}
\fi
\usepackage[english,french]{babel}
\usepackage[a4paper,margin=15mm,top=17mm,bottom=17mm]{geometry}
\usepackage{xcolor}
\usepackage{array}
\usepackage{tabularx}
\usepackage{booktabs}
\usepackage{microtype}
\usepackage{enumitem}
\usepackage{fancyhdr}
\usepackage[hidelinks]{hyperref}

\definecolor{ink}   {HTML}{1A1A1A}
\definecolor{accent}{HTML}{1F5C8B}
\definecolor{good}  {HTML}{2E7D32}
\definecolor{bad}   {HTML}{C62828}
\definecolor{warn}  {HTML}{EF6C00}
\definecolor{muted} {HTML}{6E6E6E}
\definecolor{rule}  {HTML}{D8D8D4}
\definecolor{wash}  {HTML}{F4F6F8}

\renewcommand{\familydefault}{\sfdefault}
\color{ink}
\setlength{\parindent}{0pt}
\setlength{\parskip}{0pt}

\pagestyle{fancy}
\fancyhf{}
\renewcommand{\headrulewidth}{0pt}
\fancyfoot[C]{\footnotesize\color{muted}\thepage}
\fancyhead[L]{\footnotesize\color{muted}Noise Modelling Hanoi --- speaker script}
\makeatletter\fancyhead[R]{\footnotesize\color{muted}\@scriptname}\makeatother

% --------------------------------------------------------------- structure
% \slide{number}{short title}{minute mark}
% \slide[<on-screen fraction>]{<page>}{<title>}{<minute mark>}
%
% TWO NUMBERS, because the deck shows one and the file has the other. \#1 is the
% PDF page -- what you get by pressing right that many times from the start. The
% optional argument is the FRACTION PRINTED ON THE SLIDE ITSELF, which is lower,
% because metropolis does not count section pages or the appendix. Presenters
% navigate by what they can see, so the visible number is given first and in
% bold; the page number follows it in grey for anyone opening the PDF.
\newcommand{\slide}[4][]{%
  \par\vspace{5mm}\noindent
  \colorbox{accent}{\parbox{\dimexpr\textwidth-2\fboxsep}{%
    \color{white}\vspace{.6mm}
    \if\relax\detokenize{#1}\relax
      \textbf{#2}%
    \else
      \textbf{#1}\,{\footnotesize\color{white!70}(p.\,#2)}%
    \fi
    \hspace{3mm}#3\hfill\small #4\vspace{.6mm}}}
  \par\vspace{1.5mm}}

% \section-level marker for an act of the talk
\newcommand{\act}[2]{%
  \par\vspace{7mm}\noindent
  {\color{accent}\rule{\textwidth}{.9pt}}\par\vspace{1.5mm}
  {\large\bfseries\color{accent}#1}\hfill{\color{muted}#2}
  \par\vspace{.5mm}{\color{rule}\rule{\textwidth}{.4pt}}\par}

% The two columns. \say is the whole unit: English left, French right.
%
% Fixed widths, not tabularx: tabularx scans its body for a literal
% \end{tabularx}, so wrapping it in a \newenvironment hides the terminator and
% the whole document ends up inside one unfinished cell. 84 + 5 + rule + 5 + 84
% is \textwidth on A4 with 15 mm margins.
\newcolumntype{S}{>{\setlength{\parskip}{1.1mm}}p{84mm}}
\newenvironment{say}{%
  \begin{tabular}{@{}S@{\hspace{5mm}}|@{\hspace{5mm}}S@{}}
}{%
  \end{tabular}\par}

% The deck's own shorthands. Frame titles are copied into this file VERBATIM by
% check_script.py, so anything main.tex can put in a title has to resolve here
% too -- \dB in "a $-5.1\dB$ discontinuity" is what first broke the build.
% Keep this list in step with the shorthands section of main.tex.
\newcommand{\dB}{\ensuremath{\,\mathrm{dB}}}  % \ensuremath: text here, math in the deck
\newcommand{\rr}{$R^2$}
\newcommand{\LA}{$L_{A,25\mathrm{s}}$}
\newcommand{\up}[1]{\textcolor{good}{\textbf{#1}}}
\newcommand{\down}[1]{\textcolor{bad}{\textbf{#1}}}
\newcommand{\wn}[1]{\textcolor{warn}{\textbf{#1}}}

% A line that must be said exactly, or must never be said.
\newcommand{\must}[1]{{\color{good}\textbf{#1}}}
\newcommand{\never}[1]{{\color{bad}\textbf{#1}}}
\newcommand{\care}[1]{{\color{warn}\textbf{#1}}}
\newcommand{\cue}[1]{\par\vspace{1mm}{\footnotesize\color{muted}\textit{#1}}\par}

% A boxed aside spanning the width: for questions to expect.
\newcommand{\aside}[2]{%
  \par\vspace{1.5mm}\noindent
  \colorbox{wash}{\parbox{\dimexpr\textwidth-2\fboxsep}{\vspace{1mm}
    \footnotesize\textbf{#1}\par\vspace{.7mm}#2\vspace{1mm}}}\par}

